\section{DSPy 编译与调试实践}

\subsection{日志与时间线}

DSPy 的 optimizer 产生多组 prompt 组合，建议将其划分为三个阶段的日志流：

\begin{enumerate}
    \item \textbf{输入阶段}：记录 signature、few-shot 样本与期望 metric
    \item \textbf{编译阶段}：保存 optimizer 生成的候选 prompt 与评分
    \item \textbf{运行阶段}：对比各候选在真实数据上的表现，记录 trace 与 response
\end{enumerate}

\subsection{调试小技巧}

调试 DSPy 编译器可以借助以下方法：

\begin{itemize}
    \item 将每次 optimizer 输出格式化为 human-readable YAML 并注释效果
    \item 插入 fallback 流程：若某模块输出异常则自动回退到上一次成功 pipeline
    \item 在 Evidence Matrix 中保留 prompt → response → metric 的对应关系，便于快速回溯
\end{itemize}

\begin{knowledgebox}{证据矩阵即调试日志}
Evidence Matrix 展示候选 prompt、生成答案与评分，越多的历史版本被记录，越容易定位哪一步引发偏差。把这张矩阵截图 (如[cover.jpg](/home/v-haoqiwang/ai-course-notes/talks/berkeley-llm-agents/f24/lecture08/cover.jpg)) 引入笔记，即使不是运行中帧，也传达“编译可观察”的理念。
\end{knowledgebox}

\subsection{Evidence Frame}

\begin{figure}[H]
\centering
\includegraphics[width=0.85\textwidth]{cover.jpg}
\caption{Evidence Matrix 示例： prompt → 结果 → metric 的记录表}
\end{figure}

\subsection{本章小结}

DSPy 编译需要在优化过程中记录每个候选并可重放，Evidence Matrix 是调试与回滚的第一手资料。

